\documentclass[12pt]{article}
\usepackage[utf8]{inputenc}
\usepackage[T1]{fontenc}
\usepackage{hyperref}
\title{Upper Tanana Language Resources for K-12 Curriculum Development}
\author{Compiled by Claude Sonnet 5 (claude-sonnet-5) for Andr\'as Kornai}
\date{}
\begin{document}
\maketitle

\section{Introduction}
Upper Tanana (ISO 639-3 \texttt{tau}, Glottocode uppe1437) is an Athabaskan language of eastern Alaska and the Yukon, with roughly 99 speakers according to the reference description. Sources found in one web search disagree on speaker numbers (Wikipedia mentions two known speakers at Nabesna and about 20 at Northway), and the language is classed as critically endangered. This report rests on one search and is incomplete.

\section{Primary Educational Resources}
\begin{itemize}
\item Tetlin dialect junior dictionary by Paul Milanowski and Alfred John (1979, National Bilingual Materials Development Center, University of Alaska Anchorage), made for elementary students. The first edition ran from English to Upper Tanana. Later editions (2007, 2009, 2013) are described in the UBC pages below. \url{https://knowledgebase.arts.ubc.ca/dictionary-for-learners-tetlin-dialect-of-upper-tanana-athabascan}
\end{itemize}
No curriculum, school programme, or classroom materials beyond the dictionaries were found.

\section{Linguistic and Reference Materials}
\begin{itemize}
\item UBC Relational Lexicography knowledgebase entry for the Neeaaneegn' Upper Tanana (Tetlin) dictionary. \url{https://knowledgebase.arts.ubc.ca/neeaaneegn-upper-tanana-tetlin-dictionary/}
\item Further UBC knowledgebase entries on Upper Tanana dictionary editions. \url{https://knowledgebase.arts.ubc.ca/?p=8431}, \url{https://knowledgebase.arts.ubc.ca/?p=6925}, \url{https://knowledgebase.arts.ubc.ca/?p=8434}
\item Wikipedia, ``Upper Tanana language'' (overview, secondary source). \url{https://en.wikipedia.org/wiki/Upper_Tanana_language}
\item Numerals data for Upper Tanana, Max Planck EVA. \url{https://lingweb.eva.mpg.de/channumerals/Tanana-Upper.htm}
\end{itemize}

\section{Religious and Community Resources}
The 2013 dictionary edition is listed under the Native Village of Tetlin as publisher, per the search summary. No other community organisation pages were found.

\section{Digital and Audio Resources}
\begin{itemize}
\item Endangered Languages Project entry for Upper Tanana. \url{https://www.endangeredlanguages.com/lang/1165}
\end{itemize}
UBC reports that PDF scans of the dictionary editions are open access through the Alaska Native Language Archive. I did not see that archive's URL in the results, so I give none. No audio or video resources were found.

\section{Curriculum Development Context}
The language is spoken as a first language by adults only and is not known to be taught in schools. The Tetlin dictionary, originally designed for elementary pupils, is the most useful found material and a natural base for early-grade vocabulary work. Any curriculum should be developed with the Tetlin and Northway communities and the Alaska Native Language Center and Archive, whose holdings I could not inspect.

\end{document}